% generated by report/tables/make_tables.py from avalon/llm_run.py and avalon/followup.py
\paragraph{System prompt for direct answers and thinking (rule card + task)}
\begin{Verbatim}[breaklines,breakanywhere,fontsize=\scriptsize,frame=single,rulecolor=\color{black!25},framesep=4pt]
你在分析一局《阿瓦隆》（The Resistance: Avalon）。只能依据下列规则和给出的记录做严格的逻辑推理；发言、投票倾向、"像不像"之类的策略性推测都不算证据。

【阵营与人数】好人阵营：梅林、派西维尔、亚瑟的忠臣。坏人阵营：莫甘娜、莫德雷德、奥伯伦、莫德雷德的爪牙。坏人人数：5人局2名，6人局2名，7人局3名，8人局3名，9人局3名，10人局4名。本局使用的角色在记录开头给出，每人一张。
【开局信息】每个人知道自己的角色。除奥伯伦外的坏人互相认识，他们看不到奥伯伦，奥伯伦也看不到他们。梅林能看到除莫德雷德以外的所有坏人（包括奥伯伦）。派西维尔能看到梅林和莫甘娜，但分不清谁是谁。亚瑟的忠臣和奥伯伦开局得不到任何信息。
【组队与投票】队长提名队伍，所有人公开投票，赞成票严格过半才通过（平票算否决）。每次提案后队长按座位顺序交给下一位。同一轮连续5次否决则坏人直接获胜。投票可以随意投，规则上与身份无关。各轮出队人数：5人局2,3,2,3,3；6人局2,3,4,3,4；7人局2,3,3,4,4；8至10人局3,4,4,5,5。
【执行任务】每名队员出一张任务牌：好人必须出成功；坏人可以出成功，也可以出失败。牌洗混后翻开，所有人知道失败票的张数，但不知道是谁出的。出现1张失败票任务就失败；例外：7人及以上的第4个任务需要至少2张失败票才失败。
【结束】三个任务失败则坏人获胜。三个任务成功后刺客刺杀一人：刺中梅林则坏人获胜，否则好人获胜。
【湖中仙女】持有者查验一名玩家，得知其阵营（好人或坏人），结果只有持有者本人知道。

判断给定"说法"相对于规则和记录（以及你的视角信息）的地位，只能回答一个词：
成立：在所有符合规则与记录的身份分配中，说法都为真；
不成立：在所有这样的身份分配中，说法都为假；
不确定：有的身份分配下为真，有的为假；
矛盾：记录或你的视角信息本身违反规则，不存在任何符合的身份分配。
\end{Verbatim}

\paragraph{Answer instruction appended to every question}
\begin{Verbatim}[breaklines,breakanywhere,fontsize=\scriptsize,frame=single,rulecolor=\color{black!25},framesep=4pt]
只输出 成立、不成立、不确定、矛盾 中的一个词，不要解释。
\end{Verbatim}

\paragraph{Extraction system prompt (followed by one worked example)}
\begin{Verbatim}[breaklines,breakanywhere,fontsize=\scriptsize,frame=single,rulecolor=\color{black!25},framesep=4pt]
把一段《阿瓦隆》局面描述逐字转换成JSON，不做任何推理。格式：
{"n": 人数,
 "config": {角色英文名: 数量, ...}（只列出现的角色）,
 "quests": [{"q": 任务编号, "team": [座位号], "fails": 失败票张数, "result": "success"或"fail"}]（只列已执行的任务；提案和投票不需要）,
 "view": null 或 {"player": 座位号, "role": 角色英文名, "seen": [开局看到的座位号], "lady": [[被查验的座位号, 是坏人则true否则false]]},
 "claim": 说法的JSON}
说法的写法：X号一定是坏人 → {"t":"evil","p":X}；X号一定是好人 → {"t":"good","p":X}；X号一定是梅林 → {"t":"role","p":X,"r":"merlin"}；X号一定不是梅林 → {"t":"not","a":{"t":"role","p":X,"r":"merlin"}}；某些人中至少有k名坏人 → {"t":"atleast","ps":[座位号],"k":k}；某些人中恰好有k名坏人 → {"t":"exactly","ps":[座位号],"k":k}；X号和Y号一定属于同一阵营 → {"t":"same","ps":[X,Y]}；本局的坏人恰好是某些人 → {"t":"evil_set","ps":[座位号]}。
角色英文名：梅林 merlin，派西维尔 percival，亚瑟的忠臣 servant，莫甘娜 morgana，莫德雷德 mordred，奥伯伦 oberon，莫德雷德的爪牙 minion。座位号使用文字中的号码（从1开始）。旁观者视角时 view 为 null。只输出一个JSON对象，不要解释，不要Markdown。
\end{Verbatim}

\paragraph{The worked extraction example: input}
\begin{Verbatim}[breaklines,breakanywhere,fontsize=\scriptsize,frame=single,rulecolor=\color{black!25},framesep=4pt]
6人局，角色为：梅林、派西维尔、亚瑟的忠臣×2、莫甘娜、莫德雷德的爪牙（共2名坏人）。
第1轮第1次提案：队长3号提名1号、6号，赞成者为1号、3号、6号，被否决。
第1轮第2次提案：队长4号提名2号、4号，赞成者为1号、2号、4号、5号，通过。
第1个任务由2号、4号执行，翻开1张失败票，任务失败。
你是5号，身份是派西维尔。开局时你看到1号、3号（其中一人是梅林，另一人是莫甘娜，你分不清）。
说法：2号和4号中至少有1名坏人。
\end{Verbatim}

\paragraph{The worked extraction example: output}
\begin{Verbatim}[breaklines,breakanywhere,fontsize=\scriptsize,frame=single,rulecolor=\color{black!25},framesep=4pt]
{"n": 6, "config": {"merlin": 1, "percival": 1, "servant": 2, "morgana": 1, "minion": 1}, "quests": [{"q": 1, "team": [2, 4], "fails": 1, "result": "fail"}], "view": {"player": 5, "role": "percival", "seen": [1, 3], "lady": []}, "claim": {"t": "atleast", "ps": [2, 4], "k": 1}}
\end{Verbatim}

\paragraph{Reveal-set question}
\begin{Verbatim}[breaklines,breakanywhere,fontsize=\scriptsize,frame=single,rulecolor=\color{black!25},framesep=4pt]
现在可以用湖中仙女查验若干玩家的阵营（每查验一人，就得知他是好人还是坏人）。请一次性给出最少需要查验的玩家，使得无论查验结果如何，都一定能判断上述说法成立还是不成立。不能查验你自己。只输出座位号的JSON数组，例如 [3, 5]，不要解释。
\end{Verbatim}

\paragraph{Budget-forcing suffix (appended to an unfinished thinking trace)}
\begin{Verbatim}[breaklines,breakanywhere,fontsize=\scriptsize,frame=single,rulecolor=\color{black!25},framesep=4pt]
考虑到思考时间有限，我现在必须根据已有的思考直接给出最终答案。
</think>
\end{Verbatim}

\paragraph{Extraction-repair instruction (after the solver-reported conflict)}
\begin{Verbatim}[breaklines,breakanywhere,fontsize=\scriptsize,frame=single,rulecolor=\color{black!25},framesep=4pt]
请逐项对照原文检查你的转写：人数、角色配置、每个任务的队员、失败票张数和结果、你的视角信息、说法。如果转写有误，输出改正后的完整JSON；如果转写与原文一致（也就是原文记录本身违反规则），就原样输出同一个JSON。只输出一个JSON对象，不要解释，不要Markdown。
\end{Verbatim}
